% Pour toute question sur l'utilisation de cette classe, écrivez à jerome.soucy(at)ulaval.ca
% Nécessite la classe Evaluation.sty dans le même dossier ou dans le PATH de votre système
% Nécessite le fichier logo_ul.pdf dans le même dossier un dans le PATH de votre système
\def\Ponderation{10}
\def\SigleNumero{STT-2100}
\def\NomCours{Régression}
%\def\DateEvaluation{2 mars 2023}
\def\DateRemise{6 octobre 2024}
\def\HeureRemise{23h59}
\def\HeureEvaluation{}
\def\Session{Été 2021}
\def\NomEvaluation{Devoir 1}

% Ne rien mettre entre les accolades si l'enseignant (ou l'enseignante) est aussi le coordonnateur (ou la coordonnatrice)
\def\Coordonnatrice{}
\def\Coordonnateur{}

% Ne rien mettre entre les accolades s'il y a plus d'une section
\def\Enseignant{}
\def\Enseignante{Audrey-Anne Vallée}

% Ne rien mettre entre les accolades s'il s'agit d'un cours à section unique
% Mettre le nom de chacun des enseignants et enseignantes entre les accolades
% Une case à cocher sera générée pour chacune des sections
\def\SectionA{}
\def\SectionB{}
\def\SectionC{}
\def\SectionD{}


\def\Directives{
\begin{enumerate}
\item Le devoir compte \numquestions ~questions pour un total de \numpoints{}~points.
\item À faire individuellement ou en équipe de deux personnes. 

\item Vous devez remettre \underline{deux fichiers} par équipe:
\begin{itemize}
\item Un fichier PDF avec vos réponses aux questions 1 à 5. Vous pouvez rédiger à l'ordinateur dans \LaTeX, à la main sur un papier que vous numérisez ou sur une tablette. Ce fichier est à remettre sur Gradescope, à l'aide des instructions données au point 4. sur cette page.
\item Un script \texttt{R} ou \texttt{SAS} avec les codes utilisés pour la question 5. À remettre dans la boîte de dépôt du site du cours.
\end{itemize}

\item Remises sur \textit{Gradescope}:
\begin{itemize}
\item Rendez-vous au www.gradescope.com;
\item Connectez-vous en utilisant votre adresse courriel de l'Université (idul@ulaval.ca) ainsi que votre mot de passe habituel. Il faut faire «Log in», utiliser «School credentials» et chercher Université Laval dans la liste . Vous devriez déjà avoir un compte, par le biais de l’université. 
\item Une fois sur le site, choisissez STT-2100;
\item En sélectionnant le Devoir 2, vous ajoutez le PDF qui contient vos réponses;
\item Vous devez identifier les différentes sous-questions sur les pages de votre devoir.
\end{itemize}

\item Il est \textit{\textbf{nécessaire}} de justifier toutes les démarches. Une réponse qui n'a pas une démarche \textit{\textbf{complète}} ne pourra avoir tous les points. Vous ne pouvez pas utiliser les résultats du cours directement, à moins que ce soit indiqué (par exemple, à la question 1).
\item La qualité de la langue sera évaluée, incluant l'orthographe, la grammaire et la syntaxe. Un maximum de 10 points pourra être retranché pour mauvaise qualité du français. 
\item Les travaux remis en retard recevront la note 0.
%\item Veuillez relire le point 7, il n'y aura aucune exception.
\end{enumerate}
}

% Veuillez indiquer le type d'évaluation entre crochets
% Examen : option à choisir pour un examen présentiel qui sera corrigé de manière classique. Un tableau des points sera affiché sur la page couverture.
% Gradescope : option à choisir pour un examen qui sera corrigé via Gradescope. Aucun tableau de pointage n'apparaîtra et il y a quelques différences au niveau de la mise en page
% TPS :
% TPF :
% Devoir :
\documentclass[Devoir]{Evaluation}
\usepackage{multicol,graphicx}
\begin{document}

\fbox{%
    \parbox{\textwidth}{%

Pour certains exercices, nous référerons au modèle de régression linéaire simple donné par

\begin{equation}\label{mod:reglin}
    Y_i = \beta_0\,+\,\beta_1\, x_i \,+\, \varepsilon_i\qquad \text{ pour } i=1,\ldots,n.
\end{equation}

Dans ce modèle, les 3 postulats sont respectés et nous avons $\varepsilon_i \overset{iid}{\sim} \mathcal{N}(0, \sigma^2), \text{ pour } i=1,\ldots,n$.
    }%
}



\vspace{1cm}
\begin{questions}  

\vspace{1cm}
\question[10]

Considérons le modèle de régression donné par l'équation \eqref{mod:reglin} au début du document. On a déjà déterminé la distribution de $\widehat{\beta}_0$ et celle de $\widehat{\beta}_1$. On s'intéresse à l'estimation de $\sigma^2$. \\

Montrer que 
$$\hat\sigma_{MV}^2=\frac{1}{n}\sum\limits_{i=1}^n \hat\varepsilon_i^2$$ 
est un estimateur  biaisé de la variance autour de la droite $\sigma^2$. (Vous pouvez utiliser directement  la distribution des résidus abordée dans le chapitre 4.)


\begin{solution}
    \begin{align*}
            \mathrm{E}\left(\widehat{\sigma}_{MV}^2\right)
            &=\mathrm{E}\left(\frac{1}{n} \sum\limits_{i=1}^n \hat\varepsilon_i^2\right) \cr
            &= \frac{1}{n} \sum\limits_{i=1}^n \mathrm{E}\left(\hat\varepsilon_i^2\right) \cr
             &= \frac{1}{n} \sum\limits_{i=1}^n \left[ \mathrm{Var}\left(\hat\varepsilon_i\right) + \mathrm{E}\left(\hat\varepsilon_i\right)^2 \right]\cr
             &= \frac{1}{n} \sum\limits_{i=1}^n \left[ \left(1- \frac{1}{n} - \frac{(x_i - \overline{x})^2}{S_{xx}}\right)\sigma^2+ 0\right]\cr
             &= \frac{\sigma^2}{n} \sum\limits_{i=1}^n \left[ 1- \frac{1}{n} - \frac{(x_i - \overline{x})^2}{S_{xx}}\right]\cr
             &= \frac{\sigma^2}{n}  \left[ n- 1 - 1\right]\cr
             &= \sigma^2 \frac{n-2}{n}
        \end{align*}
\end{solution}

\vspace{2.5cm}
\question

Soit le modèle de régression donné par l'équation \eqref{mod:reglin} au début du document. Nous tentons d'estimer $\beta_0$ et $\beta_1$ dans ce modèle à l'aide de $\widehat{\beta}_0^*$ et $\widehat{\beta}_1^*$ que nous définissons comme les valeurs qui minimisent la fonction:

\begin{equation*}
   Q(\beta_0, \beta_1)= \sum_{i=1}^n(Y_i-\beta_0- \beta_1x_i)^2 + \lambda \beta_1 ^2,
\end{equation*}

où $\lambda \in \mathbb{R}_+$. (Aux sous-questions (b) et (d), vous pouvez utiliser directement  les résultats du chapitre 3.)
\begin{parts}
    \part[8] Donnez les expressions de $\widehat{\beta}_0^*$ et de $\widehat{\beta}_1^*$.
    \part[4] Calculez le biais de $\widehat{\beta}_1^*$.
    \part[1] Pour quelle(s) valeur(s) de $\lambda$ a-t-on que $\widehat{\beta}_1^*$ est sans biais pour $\beta_1$?
    \part[7] Pour quelle(s) valeur(s) de $\lambda$ a-t-on $\text{Var}\left(\widehat{\beta}_1^*\right)<\text{Var}\left(\widehat{\beta}_1\right)$, où $\widehat{\beta}_1$ est l'estimateur des moindres carrés ordinaires?
    \part[4] Quel serait l'intérêt d'utiliser une valeur de $\lambda$ qui induit un biais dans notre estimation de $\beta_1$?
    
\end{parts}

\begin{solution}
    \begin{parts}
        \part \begin{align*}
            \frac{\partial Q}{\partial \beta_0}
            &= -2 \sum (Y_i-\beta_0- \beta_1x_i) \cr
            %
            \frac{\partial Q}{\partial \beta_1}
            &= -2 \sum (Y_i-\beta_0- \beta_1x_i)x_i + 2\lambda \beta_1 
        \end{align} 
        \begin{align*}
            \sum (Y_i-\widehat{\beta}_0^* - \widehat{\beta}_1^* x_i) 
            &= 0 \cr
            n\overline{Y} - n\widehat{\beta}_0^* - n\widehat{\beta}_1^* \overline{x}
            &= 0 \cr
            \overline{Y} - \widehat{\beta}_1^* \overline{x}
            &=\widehat{\beta}_0^* \cr
            %
            \sum (Y_i-\widehat{\beta}_0^* - \widehat{\beta}_1^* x_i) x_i -\lambda \widehat{\beta}_1^* 
            &= 0 \cr
            \sum Y_ix_i - n\widehat{\beta}_0^*\overline{x} - \widehat{\beta}_1^* \sum x_i^2 -\lambda \widehat{\beta}_1^*
            &= 0 \cr
            \sum Y_ix_i - n(\overline{Y} - \widehat{\beta}_1^* \overline{x})\overline{x} - \widehat{\beta}_1^* \sum x_i^2 -\lambda \widehat{\beta}_1^* 
            &= 0 \cr
            \sum Y_ix_i - n\overline{Y}\overline{x} + n\widehat{\beta}_1^* \overline{x}^2 - \widehat{\beta}_1^* \sum x_i^2 -\lambda \widehat{\beta}_1^* 
            &= 0 \cr
            S_{xy}
            &=  \widehat{\beta}_1^* S_{xx} +\lambda \widehat{\beta}_1^* \cr
            \frac{S_{xy}}{S_{xx} + \lambda}
            &=  \widehat{\beta}_1^* *
            \cr
        \end{align} 
        \part  $$ \mathrm{E} (S_{xy}/S_{xx}) = \beta_1  \rightarrow \mathrm{E} (S_{xy}) = \beta_1 S_{xx}  $$
        $$ \mathrm{E} (S_{xy}/(S_{xx}+\lambda)) = \mathrm{E} (S_{xy})/(S_{xx}+\lambda) = \beta_1 S_{xx} / (S_{xx}+\lambda) $$
        \part $\lambda = 0$
        \part $$ \mathrm{Var} (S_{xy}/S_{xx}) = \sigma^2/S_{xx}  \rightarrow (1/S_{xx}^2)\mathrm{Var} (S_{xy}) = \sigma^2/S_{xx}  \rightarrow \mathrm{Var} (S_{xy}) = \sigma^2 S_{xx}  $$
         $$ \mathrm{Var} (S_{xy}/(S_{xx}+\lambda)) = \mathrm{V} (S_{xy})/(S_{xx}+\lambda)^2 = \sigma^2 S_{xx}  / (S_{xx}+\lambda)^2 $$
         $$  \frac{\sigma^2 S_{xx}}{(S_{xx}+\lambda)^2} \leq \frac{\sigma^2}{S_{xx}} \text{si}  
         \frac{ S_{xx}^2}{(S_{xx}+\lambda)^2} \leq 1 \text{si}  
         S_{xx}^2 \leq (S_{xx}+\lambda)^2 \text{si}  
         S_{xx}^2 \leq S_{xx}^2 + 2S_{xx}\lambda +\lambda^2 
         \text{si}  
         0 \leq  2S_{xx}\lambda +\lambda^2$$
         pour $\lambda > 0.$
         \part Diminuer l'EQM, avoir une plus petite variance.
    \end{parts}
\end{solution}


\newpage
\question 
Soit le modèle de régression donné par l'équation \eqref{mod:reglin} au début du document. Nous tentons d'estimer $\beta_0$ et $\beta_1$ en nous limitant à la classe des estimateurs linéaires.
Les estimateurs $\widehat{\widetilde{\beta}}_0$ et $\widehat{\widetilde{\beta}}_1$ sont linéaires s'ils sont de la forme 
\begin{align*}
    \widehat{\widetilde{\beta}}_0=\sum_{i=1}^n  \alpha_i Y_i, & & \widehat{\widetilde{\beta}}_1=\sum_{i=1}^n  \lambda_i Y_i,
\end{align*}
où $\alpha_1,\hdots,\alpha_n, \lambda_1,\hdots, \lambda_n$ sont des constantes connues.

\begin{parts}
    \part[4] Montrez que $\widehat{\widetilde{\beta}}_0$ est sans biais pour $\beta_0$ si les deux conditions suivantes sont respectées:
    \begin{align*}
        \sum_{i=1}^n  \alpha_i =1, &   & \sum_{i=1}^n  \alpha_i x_i =0.
    \end{align*}
    \part[4] Quelles conditions pouvons-nous imposer aux $\lambda_i$ pour que $\widehat{\widetilde{\beta}}_1$ soit sans biais pour $\beta_1$?
    \part[4] Montrez que l'estimateur linéaire non-biaisé de $\beta_1$ ayant la plus petite variance parmi tous les estimateurs linéaires non-biaisés de $\beta_1$ est celui qui minimise la valeur de $\displaystyle\sum_{i=1}^n\lambda_i ^2$.
    \part[2] Donnez l'expression de $ \alpha_i$, dans le cas de l'estimateur des moindres carrés pour $\beta_0$.
\end{parts}


\begin{solution}
    \begin{parts}
        \part $$\mathrm{E} = \sum \lambda_i (\beta_0 + \beta_1 x_i)$$
        \part $$\mathrm{E} = \sum \alpha_i (\beta_0 + \beta_1 x_i) = \beta_0 \sum \alpha_i + \beta_1 \sum \alpha_i x_i$$
        $ \sum \alpha_i = 1$ et $\sum \alpha_i x_i = 0$
        \part $$ \mathrm{Var}  = \sum \lambda_i^2 \sigma^2$$
        \part $\widehat{\beta}_0 = \overline{Y} - \widehat{\beta}_1 \overline{x}$, les $d_i$ du cours.
    \end{parts}
\end{solution}


\vspace{2.5cm}
\question 
Soit le modèle de régression donné par l'équation \eqref{mod:reglin} au début du document. 
Nous supposons que les $x_i$ sont des variables binaires, ne pouvant prendre que les valeurs $0$ ou $1$.
Nous notons
$$  n_0=\sum_{i=1}^n (1-x_i); \quad
    n_1=\sum_{i=1}^n x_i; \quad
    \overline{Y}_0=\frac{1}{n_0}\sum_{i=1}^n (1-x_i)Y_i; \quad
    \overline{Y}_1=\frac{1}{n_1}\sum_{i=1}^n x_iY_i.$$


\begin{parts}
    \part[4] Dans vos mots, dites ce que représentent les quantités $n_0$, $n_1$, $\overline{Y}_0$ et $\overline{Y}_0$.
    \part[4] Montrez que l'estimateur des moindres carrés pour $\beta_0$ est donné par  $\widehat{\beta}_0=\overline{Y}_0$. Vous pouvez utiliser le résultat $\widehat{\beta}_0= \overline{Y}-\widehat{\beta}_1\overline{x}$.
    \part[4] Montrez que l'estimateur des moindres carrés pour $\beta_1$ est donné par  $\widehat{\beta}_1=\overline{Y}_1-\overline{Y}_0 $. Vous pouvez utiliser le résultat $\widehat{\beta}_1= \displaystyle\frac{S_{xY}}{S_{xx}}$.
\end{parts}

\begin{solution}
    \begin{parts}
        \part \begin{itemize}
            \item $n_0$: le nombre d'individus ayant $x=0$.
            \item $n_1$: le nombre d'individus ayant $x=1$.
            \item $\overline{Y}_0$: la moyenne de $Y$ des individus ayant $x=0$.
            \item $\overline{Y}_1$: la moyenne de $Y$ des individus ayant $x=1$.
        \end{itemize}
        \part 
        \begin{align*}
            \widehat{\beta}_0 &= \overline{Y}-\widehat{\beta}_1\overline{x} \cr
            &= 
        \end{align*}
    \end{parts}
\end{solution}

\newpage

\question\label{prog}
Les tuyaux d'égouts municipaux sont soumis à la corrosion et à des bris d'usure.  On s'intéresse à la taille des dommages à la paroi extérieure des égouts. Il est possible de mesurer la taille d'un dommage sur le terrain (sans creuser ni entrer dans l'égout) à l'aide d'un appareil de mesure à ultrasons. On obtiendrait une meilleure mesure en déterrant le tuyau, mais cela implique des opérations coûteuses. %On s'est soumis à l'exercice pour un certain nombre de tuyaux, en espérant obtenir un modèle stochastique prédictif pour l'avenir.
Un modèle linéaire pourrait permettre de prédire la taille véritable d'un dommage  à partir de la mesure ultrasons et ainsi évaluer si un tuyau est très endommagé sans le déterrer. \\

Chaque équipe recevra un jeu de données de 90 mesures de dommages prises sur des tuyaux différents. %, sélectionné au hasard parmi un plus grand jeu de données.
Pour chaque tuyau, on vous donne «\texttt{Terrain}», la taille (en mm) de la détérioration mesurée par ultrasons, et «\texttt{Reel}», la taille (en mm) mesurée en déterrant le tuyau.  Vous devez rédiger des parties d'un rapport d'analyse de régression en répondant aux questions suivantes. Utilisez $\alpha = 0,01$ pour toutes vos inférences. \\

{\bf Dans le fichier PDF}: identifiez clairement la question associée à votre réponse (par exemple (i) a)). Donnez les réponses aux questions et les tableaux de résultats obtenus avec votre logiciel qui sont nécessaires pour
répondre aux questions. Les figures et tableaux inclus dans vos réponses doivent être présentés de façon professionnelle. Ils doivent avoir des titres clairs et pertinents. Chaque figure présentée doit être discutée dans vos réponses. 

\textbf{Sorties statistiques:}
Vous devez fournir les programmes R ou SAS utilisés pour obtenir vos réponses. Un programme qui ne fonctionne pas ou qui est mal annoté entraînera une pénalité. \\


\begin{parts}
	\part[9] {\bf Description du jeu de données et de l'objectif poursuivi.}
\begin{enumerate}[(i)]
\item Identifiez et définissez les variables qui nous intéressent.
%\item Mentionner le nombre et la source des observations.
\item Décrivez {\it brièvement} la distribution de chacune des variables (graphiques et statistiques descriptives). Commentez la distribution conjointe des variables à l'aide d'un diagramme de dispersion.  (8 lignes maximum, sans compter les figures).
\item Décrire l'objectif de l'analyse de régression. (4 lignes maximum).
\end{enumerate}
	\part[14] {\bf Choix de modèle utilisé.}
\begin{enumerate}[(i)]
	%\item Si on choisit de transformer les variables avant de procéder \`a l'analyse, justifier son choix (raisons théoriques, expérience antérieure avec ce type de données, intuition seulement\,?).
	%\item Au besoin, simplifier le mod\`ele (par exemple, en refaisant la régression sans $\beta_0$).
	\item Vérifiez l'adéquation du mod\`ele RLS et le respect des postulats au moyen de graphiques et tests pertinents. Modifiez le modèle de façon à ce que vous soyez satisfait de la justesse des postulats. (On ne peut pas toujours modifier les variables pour que tous les postulats soient exactement vérifiés;  choisissez alors le meilleur modèle possible). Dans le fichier réponse, formulez clairement le modèle et les postulats du mod\`ele retenu (i.e. écrire le modèle complet en language mathématique). Présentez seulement le modèle final, et non tous les modèles ajustés. Motivez le choix de modèle final (8 lignes maximum). Incluez toutes les étapes de votre démarche (différents modèles, vérification des postulats) dans votre programme R ou SAS. 
\end{enumerate}
	\part[17] {\bf Présentation des résultats.}
\begin{enumerate}[(i)]
	\item Donnez les estimations des param\`etres (${\hat \beta}_0$, ${\hat \beta}_1$, ${\hat \sigma}^2$).
	\item Fournissez des intervalles de confiance pour les paramètres (${\hat \beta}_0$, ${\hat \beta}_1$) et une interprétation en contexte. (8 lignes maximum pour l'interprétation).
	\item Discutez de la signification statistique des coefficients ${\hat \beta}_0$ et ${\hat \beta}_1$ à l'aide de tests $t$. (6 lignes maximum).
	\item L'analyse de variance permet de faire un test d'hypothèses. Donnez les hypothèses de ce test et discutez de sa signification statistique pour vos données. (4 lignes maximum pour la discussion).
\end{enumerate}
\end{parts}

\end{questions}


\end{document}




\vspace{0.3in}
\textbf{Question 4 (10 pts)} (Design d'expérience)\\

On s'intéresse à un médicament sensé réduire le nombre de crises chez les individus atteints d'epilepsie. Vous devez aider à planifier une expérience pour étudier l'effet de la dose du médicament sur le nombre de crises. Puisque vous vous intéressez à la relation entre deux variables quantitatives, vous planifiez analyser les données avec un modèle de régression linéaire simple. \\

Soit $x_1, \ldots, x_n$ les valeurs de la dose pour les 20 individus de votre étude.\\ Pour chacun, vous pouvez attribuer n'importe quelle dose entre 2 mg et 21 mg. \\

\begin{itemize}
%\item [(a)] Pour le modèle de régression linéaire simple $Y_i = \beta_0 + \beta_1 x_i + \varepsilon_i$, prouvez que la variance de $\beta_0$ est minimale si les $x_i$ sont tels que $\bar{x}$ = 0.
\item [(a)] (5 pts) Quelle est la meilleure fa\c{c}on de choisir les valeurs $x_1,\ldots,x_n$ pour minimiser la variance de votre estimateur $\hat{\beta_1}$? Justifiez votre réponse.
\vspace{0.1 in}
\item [(b)] (5 pts) Pourquoi ne serait-ce pas vraiment une bonne idée de choisir les valeurs de $x_i$ de cette façon, même si cela minimise la variance de votre estimateur?
\end{itemize}

\end{document} 